This research has established a definitive link between AGN activity, their fractional contribution to the host galaxy's luminosity, and a consequential, quantifiable impact on the galaxy property diagnostics, notably within the UVJ and ugr colour spaces. Building upon these findings, the following discussion will explore the precise mechanisms by which AGN influence the evolutionary trajectories of their host galaxies' colours and address some of the limitations, uncertainties, and inherent assumptions in this work.
\section{The Lacy Wedge for AGN Identification}
Mid-infrared colours have long been used to select AGN in large photometric surveys. To that end, the IRAC colour-colour diagram was used to identify if the AGN models were effectively behaving as they did in observation. While many AGN diagnostics exist, one of the most widely used is the Lacy Wedge \citep{lacy_obscured_2004}. By utilising these diagnostics, it was clear that both Type 1 and Type 2 AGN models affect the MIR colourspace. Importantly, we see that the models do replicate the expected behaviour of a Type 2 AGN where there is an initial low completeness value, and introducing an AGN fraction drastically increases the selection completeness (See Figure \ref{fig:Brown-Type2Composite_IRAC}). This result is consistent with the design of the Lacy selection region, which was designed to select obscured AGN due to their thermally produced IR emissions from the dusty torus \citep{lacy_obscured_2004, lacy_optical_2007}. However, the Type 1 AGN composites exhibited a contrasting behaviour (See Figure \ref{fig:Brown-Type1Composite_IRAC}), clustering outside the selection region and demonstrating a significant decrease in completeness. This clustering and reduced completeness is likely attributed to the higher emissions in the shorter end of the MIR range, characteristic of Type 1 AGN \citep{thorne_deep_2022} resulting in ``bluer'' IRAC colours than those produced by a Type 2 AGN and ultimately shifting these AGNs outside the defined selection criteria.  Therefore, both Type 1 and Type 2 AGN can significantly influence AGN identification in the IRAC colour space, highlighting the effectiveness of the developed models in capturing these diverse AGN properties. \\\\
While this technique was useful for AGN selection, requiring only three photometric bands, there are certain downsides. The Lacy wedge was designed to identify AGN through their distinctive IR colours and is susceptible to contamination from star-forming galaxies exhibiting strong mid-infrared emission due to Polycyclic Aromatic Hydrocarbons (PAH; \citealt{sajina_simulating_2005}). As the models employed in this study do not explicitly account for such contamination, a direct assessment of the contamination in these colour spaces was not feasible. Future work could address this limitation by incorporating a representative sample of non-composite galaxy templates into the analysis. This would enable a more comprehensive evaluation of the Lacy wedge's efficacy in discriminating between AGN and star-forming galaxies within the IRAC colour-colour diagram. Ultimately, the IRAC colour-colour diagram was a critical tool for evaluating the effectiveness of the AGN composite models, ensuring the created composites could be identified as AGN candidates. Analysis within this diagram confirmed that the composite models successfully reproduced the expected IR colours of AGN, validating their ability to accurately represent the infrared properties of active galaxies. It is worth noting that while the Lacy Wedge was used to validate AGN candidates, it is prone to high contamination and does not effectively select Type 1 AGN \citep{donley_identifying_2012}. Future work may benefit from the use of use a different IR colour space to probe both Type 1 and Type 2 AGN effects; this could be done with either the KI and KIM methods \citep{messias_new_2012, messias_dependency_2014} or by using WISE colours \citep{caccianiga_wise_2015}. Having established the validity of the composite AGN models in the IRAC colour space, we now turn our attention to the rest-frame UVJ diagram. The UVJ offers a complementary perspective on the influence of AGN, allowing us to investigate how including an AGN component affects galaxies' classification.
\section{The Effect of AGN on the rest-frame UVJ Diagram}
The inclusion of an AGN component into the GALSEDATLAS templates caused an appreciable effect on both the galaxy fractions and the positions of sources in the UVJ (See Figure \ref{fig:Brown-Type1Composite_UVJ} and \ref{fig:Brown-Type2Composite_UVJ}). Including the Type 1 component caused a drastic drop in the quiescent fraction of galaxies, whereby quiescent selected sources moved into the star-forming region of the UVJ diagram. The cause of this is due to the increased UV radiation that is being picked up by the U band in the SED. (See Figure \ref{fig:composite_uvj}). This increase in ultraviolet radiation is characteristic of Type 1 AGN in which the central accretion disc produces short wavelength emissions in the UV and X-ray bands \citep{nour_association_2023}. As quiescent composites contain higher AGN components, they will eventually be misidentified by the UVJ diagram as star-forming sources, drastically affecting the identification of quiescent type galaxies within the colourspace due to their extremely blue colours \citep{florez_exploring_2020}.\\\\ The inclusion of an AGN component significantly altered the UVJ colours of galaxies, particularly those near the star-forming/quiescent boundary (See Figure \ref{fig:Brown-Type1Composite_UVJ}).  This was expected, as the GALSEDATLAS templates were constructed from real galaxy spectra, which inherently include contributions from intrinsic AGN activity \citep{brown_atlas_2014}. Thus, the initial position of these sources in the UVJ diagram was likely already influenced by pre-existing AGN.
Consequently, an increase in AGN activity, particularly in galaxies near quiescent or dusty star-forming regions, could amplify their UV emission. This heightened UV output can potentially shift these galaxies across the boundary and into the star-forming region of the diagram. Such migration could lead to misclassification if the contribution of AGN is not carefully considered during galaxy classification based on UVJ colours.\\\\
A detailed examination of the colour evolution of an average galaxy reveals similar behaviours between quiescent and dusty sources (See Figure \ref{fig:ZFOURGE-EAZY-UVJSingleGalaxy_Comparison}). However, it is crucial to acknowledge that the modelling approach that was used in this work does not incorporate the potential effects of dust reddening \citep{bruzual_stellar_2003}. Our models incorporate a simplification regarding dust reddening: they assume that the AGN component is introduced into a galaxy that may already be affected by dust reddening when light from a source is absorbed and scattered by dust grains along the line of sight. However, dust reddening should be applied after the AGN is incorporated. This is because the light from the combined AGN-galaxy system should be treated as a single source that gets reddened by dust. To improve the accuracy of our models, a more realistic approach would involve applying a dust reddening correction after the AGN component has been integrated into the galaxy template. This would ensure that the dust reddening accurately reflects the attenuation of the combined light from the AGN and the host galaxy.
Consequently, due to this simplification, the current representation of sources in the dusty region may not be entirely accurate. This underscores the intricate challenge of disentangling the effects of star formation, dust attenuation and AGN activity when analysing galaxy populations through UVJ diagnostics. \\\\ Unlike the Type 1 composite, the Type 2 composite did not significantly impact the colours of the sources in the UVJ diagram. In Figure \ref{fig:Brown-Type2Composite_UVJ}, there is almost no movement in the UVJ diagram, nor is there any significant change to the galaxy fractions. This was because Type 2 AGN are obscured sources that thermally reprocess much of their UV-Optical emissions into the MIR \citep{padovani_active_2017}. The UVJ filter system, encompassing ultraviolet, optical, and near-infrared wavelengths, lacks the sensitivity to detect changes induced explicitly by obscured AGN activity. Moreover, the obscuring dusty torus surrounding the AGN reprocesses the UV and optical emission, further masking any discernible variations in these wavebands and their associated derived colours. Ultimately, the impact of both Type 1 and Type 2 AGN behave consistently and predictably.
Given that Type 1 AGN can induce a `blueing' effect on their host galaxies' colours, it's crucial to investigate the potential implications for redshift determination. If not adequately accounted for, the presence of an AGN could lead to an underestimation of redshift. To explore this further, we examine how AGN influence the ugr colour space, which is particularly sensitive to the colour variations exhibited by redshifted galaxies.
\section{How AGN Impact Redshift Determination in the ugr Diagram}
The migration of sources outside the redshift selection region and the associated drop in completeness is evident in the ugr diagram of Type 1 composites (See Figure \ref{fig:Brown-Type1Composite_ugr}), with a less pronounced effect observed in Type 2 composites (See Figure \ref{fig:Brown-Type2Composite_ugr}). This highlights the significant impact of AGN, particularly Type 1 AGN, on the colour-based selection of galaxies.\\\\
Analysis of sources in redshift bins reveals that the influence of AGN emission on colour selection is strongly redshift-dependent. The target redshift selection range in the ugr diagram exploits the characteristic drop in flux in the u-band, indicative of Lyman-break galaxies \citep{giavalisco_lyman-break_2002}. However, including an AGN component introduced significant UV emission, effectively flooding this u-dropout (See Figure \ref{fig:composite_ugr}). Consequently, the impact of AGN emission is more pronounced at target redshifts (2.6 $<$ z $<$ 3.5), where the Lyman break falls within the u-band. The u-band flux is less affected by the Lyman break at lower redshifts (z $<$ 2.6), minimising the influence of AGN-induced UV emission \citep{madau_high-redshift_1996}. This explains the relative stability of low-redshift galaxies in ugr colour space, even with increasing AGN contribution.\\\\
This technique does not effectively select high-redshift (z $>$ 3.5) galaxies due to the partial dropout in the g-band and the complete dropout in the u-band. This causes these galaxies to occupy a region in the ugr diagram outside the targeted u-dropout selection window. A closer examination revealed a significant shift in the position of high-redshift galaxies, even with minimal AGN contribution. This shift is driven by the increase in u-band flux attributed to the AGN (See Figure \ref{fig:Brown-Type1Composite_ugr}). This observation is supported by the fact that sources containing AGN may appear bluer \citep{pierce_effects_2010} and thus share a location with lower redshift sources in the ugr diagram. Ultimately, both sources in the high and target redshift ranges appear to be heavily influenced by even a low contribution from an AGN, exhibiting mean vector offsets (See Figure \ref{tab:population_vectoroffset_z}) far higher than comparable contributions at low redshifts (z $<$ 2.6). This underscores the importance of considering AGN contribution when employing colour-based selection techniques for galaxies, especially at higher redshifts. \\\\ While Type 1 AGN exert a noticeable influence on the colour space, the same cannot be said for Type 2 AGN. In general, including a Type 2 AGN did not significantly impact sources within the u-dropout selection region. This disparity arises because Type 2 AGN obscures a substantial portion of their UV emission with a dusty torus \citep{hickox_obscured_2018}, effectively preventing any significant increase in u-band flux that could otherwise disrupt the dropout signature. However, a closer examination of the colour space revealed subtle movements in the upper right-hand side of the ugr diagram, a region typically associated with high-redshift galaxies \citep{giavalisco_lyman-break_2002}. This minor contamination effect could be attributed to the increased flux at longer wavelengths characteristic of Type 2 AGN.  As demonstrated by \cite{ciesla_constraining_2015}, composite SEDs of Type 2 AGN-host galaxies exhibit a greater increase in flux at the long-wavelength end. This suggests that the subtle shifts observed in the ugr diagram may be driven by the unique spectral properties of Type 2 AGN, particularly their enhanced emission at longer wavelengths, which becomes more prominent in high-redshift sources. With the potential influence of AGN on both redshift determination and galaxy classification established, we now turn to a critical comparison of our theoretical predictions with observational data. This discussion will assess the validity of our models and identify any discrepancies warrant further investigation.
\section{Comparison of Theoretical Models to Observational data}
The UVJ and ugr colours of the theoretical composite galaxy models exhibited a clear dependence on the AGN contribution. In particular, the theoretical models seemed to show that Type 1 AGN tended to make the colours of the UVJ and ugr far bluer, causing star-forming galaxies to be overestimated in the UVJ and redshift to be underestimated in the ugr. Observational data from decomposed sources in the ZFOURGE survey support these theoretical predictions.  Analysis using CIGALE (See Figure \ref{fig:uvj-cigale-comparison}) reveals that as the AGN component is removed from a galaxy, its colours shift towards the redder end. This is evidenced by an increase in the quiescent galaxy fraction and a corresponding decrease in the star-forming fraction. 
The observed increase in the quiescent galaxy fraction, coupled with a decrease in the star-forming fraction, suggests a significant impact from removing the AGN component. This reddening effect aligns with composite models where adding an AGN leads to a bluer galaxy characterised by a reduced quiescent fraction and an increased star-forming fraction. However, while the theoretical model utilises a binary classification of Type 1 and Type 2 AGN, they do not consider intermediate types where the observer's line of sight coincides with the dusty torus opening angle. This exclusion may contribute to discrepancies between the observed decomposition and theoretical models, particularly regarding the shift observed in the UVJ data. Specifically, while the UV shift aligns with theoretical predictions, the VJ shift shows a significant deviation.  As depicted in Figure \ref{fig:decomposed-shift}, this pronounced VJ shift could be attributed to the absence of intermediate-type AGN in the models.  These AGN, with their line-of-sight aligned with the dusty torus opening angle, are expected to exhibit substantial mid-infrared (MIR) emission \citep{ciesla_constraining_2015}, potentially leading to a VJ shift and subsequent misclassification as a dusty source. \\\\
Analysis of the redshift bins calculated from the observational UVJ colours (See Figure \ref{fig:redshift-comparison}) reveals an intriguing trend. Decomposed galaxy colours tend to shift towards the dusty region of the UVJ colourspace, particularly at high redshifts where the offset significantly exceeds the uncertainty for each bin. This effect becomes less pronounced at redshifts below 2, with no observable shift in the lowest redshift bin. Several factors could contribute to this redshift-dependent shift. One plausible explanation is the increased star formation activity in the early universe, peaking at a redshift of 2-3 \citep{madau_cosmic_2014}. This epoch of heightened star formation is also expected to coincide with increased black hole activity. Both processes can influence the observed colour separation, as star formation tends to produce bluer colours.
another explanation could be due to failed decomposition techniques; this could result in an overestimation of the AGN component within the source \citep{ciesla_constraining_2015}, and thus, removing this component has the potential to shift the colours. Despite these potential issues, the full and decomposed galaxy samples follow the expected evolutionary trend. Galaxies appear more star-forming at higher redshifts and become increasingly quiescent towards the present day, consistent with our current understanding of cosmic star-formation history \citep{kim_cosmic_2024}. \\\\ Both the ugr and UVJ diagrams generated from the ZFOURGE data using CIGALE (Fig. \ref{fig:ugr-diagram-decomposed}) exhibit a prominent pattern of striations, particularly noticeable when the AGN component is removed. This reveals distinct `finger-like' structures extending from the central locus. As the striations became more noticeable once the AGN component was removed, the most likely explanation is due to degeneracies in the SED parameters during the CIGALE fitting procedure \citep{yang_x-cigale_2020}. This would cause many sources to be modelled in a simple way based on the set of parameters, and thus, this would not accurately reflect the true diversity of sources that is present in ZFOURGE.
Furthermore, the ugr diagram displays a significant absence of high-redshift sources, with only one in the top right corner. This scarcity indicates a potential limitation in the SED fitting and photometric sampling procedure (Section 3).  Further investigation revealed that some high-redshift galaxies lacked adequate flux data for specific wavelength ranges, hindering their proper sampling by astSED.
Additionally, a group of high-redshift sources with an u-g value 30 was found outside the diagram's boundaries. This anomaly suggests potential inaccuracies in the photometric sampling of these galaxies' SEDs, possibly due to incorrect scaling, leading to erroneous colour calculations. Ultimately, there is still a strong agreement between the decomposed observational data and the theoretical models, which demonstrate the usefulness of these models in predicting the observed behaviour of galaxy colours within our universe using limited photometric data. While this approach for investigating the impact of AGN is relatively robust, there are certain limitations and assumptions in this work.
\section{Uncertainties and Limitations}
\subsection{Limitations in the AGN Models and Assumption of Galaxy templates}
A notable limitation in this study is the use of the AGN models and galaxy templates. This study employed AGN models with a simplified parameter selection approach to represent two AGN types.  Type 1 AGN were assumed to be observed at an inclination of 0°, while Type 2 AGN were assumed to be observed at an inclination of 90°. As discussed above, this binary classification neglects to account for intermediate-type AGN, which may represent a significant portion of AGN within the universe. Furthermore, the dust profile of the torus was modelled with a uniform distribution, characterised by an average optical depth of 7, a radial density gradient of 0.5, and a polar density gradient of 0. However, a more comprehensive range of parameter combinations warrants exploration, particularly higher optical depths that could increase infrared reprocessing and affect the colour spaces (e.g., UVJ and IRAC diagrams). A more robust approach would involve systematically investigating these parameters to identify the most appropriate or readily applicable set \citep{gonzalez-martin_role_2023}.\\\\ 
While the AGN models employed in this thesis are widely used in the literature, they rely on a simplified representation of the dusty torus, assuming a clumpy two-phase structure for reprocessing emissions from the central disk.  However, these models fail to reproduce two key features observed in AGN: the absence of silicate emission or absorption at 97000$\AA$ and the steep infrared slope between 30000 and 60000$\AA$ \citep{cruz_modeling_2023}. This discrepancy suggests that current AGN torus models may be incomplete and do not fully capture the complexities of AGN structure, emission mechanisms, and torus composition.  In particular, the SKIRTOR models used here employ a clumpy model for torus composition, which inherently influences the reprocessing of light and could lead to significant over- or underestimation of the infrared contribution. This, in turn, may have skewed the results presented in the colour-space diagrams, especially those sampling the infrared portion, such as the UVJ and IRAC diagrams.\\\\ While the GALSEDATLAS theoretical templates encompass a wide range of galaxy types, the vast diversity of galaxies in the Universe limits the completeness of any template set. Notably, the templates used here were derived from local, low-redshift galaxies \citep{brown_atlas_2014}, potentially failing to capture redshift-dependent phenomena characteristic of the early Universe, reddening due to dust, and AGN \citep{polletta_spectral_2007}. EAZY-fitted ZFOURGE templates \citep{brammer_eazy_2008, forrest_zfourge_2018} were incorporated to address this. However, while useful, these templates were derived from photometric data and fitted on a small template set (approximately 8 templates). This limited representation could hinder the accurate characterisation of the true diversity of galaxies.  It's important to acknowledge the trade-off: while more extensive template libraries can introduce degeneracies and make it difficult to distinguish between similar templates, smaller libraries may lack the breadth to encompass the full range of galaxy properties. \citep{steinhardt_templates_2023}. \\\\It is crucial to consider that the composite models were constructed by adding AGN components to galaxy templates that were assumed to have negligible intrinsic AGN activity. However, the absence of a specific flagging system for AGN hosts means that some sources with inherent AGN activity, such as Seyfert galaxies and quasars, may have been included in the initial galaxy template set. This could result in underestimating the total AGN contribution in the final colours, as the intrinsic AGN emission from these sources would not be accounted for. To address this potential bias, future iterations of the models could incorporate a pre-selection step to identify and remove confirmed AGN host galaxies from the galaxy template library. \\\\
Error estimation was not included in the original GALSEDATLAS template + AGN composites, as these models were purely theoretical. However, to enhance the robustness of future analyses, incorporating error estimation into composite models is recommended. A potential approach to achieve this involves utilising a Markov Chain Monte Carlo (MCMC) simulation. This would allow for the creation of theoretical composite models for each source, consisting of a base galaxy component and an AGN component. To derive errors in each colour, a set of composites could be generated with randomly selected parameters at a fixed inclination (to differentiate between Type 1 and Type 2 AGN). Similarly, it would also be possible to generate a suite of galaxy templates using CIGALE from multi-wavelength photometry if flux errors are readily available; a similar MCMC approach may also be taken to generate a set of galaxy templates that can be reduced to estimate errors in each of the required photometric passbands.\\\\
By introducing noise into the composite generation process for each source, a diverse range of theoretical composites can be produced, spanning the entire parameter space. This would enable the calculation of a mean SED and the estimation of standard errors based on the scatter of the parameters. Consequently, this process would yield models with associated flux uncertainties. These models and their uncertainties could then be propagated into colour-colour space for each diagram and averaged to establish a robust error metric for the colour-colour diagrams. This approach would provide a more comprehensive and statistically sound analysis of the influence of AGN on galaxy colours. Redshift scaling was not applied to the restframe templates in the initial GALSEDATLAS + AGN composites; instead, it was due to artificial redshifting. This resulted in the initial theoretical models having a uniform distribution of redshifts. This was rectified when the ZFOURGE galaxies were used. Still, a more robust approach would have been to sample the shape of the ZFOURGE distribution or a similar photometric survey to find the shape of the distribution of redshifts and then assign a redshift to each composite to create a similar redshift distribution. 

\subsection{Limitations in the SWIRE Templates} Initially, the SWIRE templates \citep{polletta_spectral_2007} were chosen as the main theoretical templates for investigating AGN contamination in colour space. The SWIRE have galaxies representing distinct morphological galaxy types such as spiral, elliptical and lenticular types and also includes other galaxy types with AGN components present, such as the Seyfert galaxies \cite{seyfert_nuclear_1943}. The initial set of composites that was constructed in the UVJ (See Figure \ref{fig:subplots}) showed an obvious trend of movement towards a bluer region. However, after investigating these colour spaces, it was determined that the SWIRE template wavelength range was inappropriate for the ugr colourspace. This was determined because the SWIRE template's short-end wavelength range begins at 1005$\AA$, a value far above the rest-frame Lyman limit of $912\AA$. The Lyman break technique used in the ugr colourspace utilises a dropout of flux at the Lyman limit in u band ($\lambda_{eff}\approx 3500\AA$) at the target redshift range\citep{giavalisco_lyman-break_2002}. Because of this, when redshifting the templates beyond a redshift of 2.5, a lack of flux is observed due to the incomplete galaxy template, leaving the SWIRE templates unusable for composite creation and photometric sampling. This truncation of the Type 1 AGN models significantly reduces their UV wavelength coverage, potentially impacting the colour analysis.  Since Type 1 AGN are believed to have the most pronounced effects on UV colours, this limitation may affect the accuracy of the derived colour transformations, particularly at higher redshifts where the UV emission is redshifted into the observed optical bands.\\
A further limitation stems from the inherent properties of the SWIRE templates. Constructed from local galaxies, these templates do not accurately represent the diversity of galaxies found across cosmic time, particularly those at redshifts of 2-3 when the cosmic star formation rate peaked \citep{dominguez_extragalactic_2011, madau_cosmic_2014}. This discrepancy limits the applicability of the SWIRE templates, especially when examining the UVJ and ugr colour evolution of galaxies hosting Type 1 AGN.